\pagebreak

\byline{{\textbf{\Large\sffamily Mini Bonsai Workshop:}}\\
        {\textbf{\large\sffamily Exploring the Art of Mame Bonsai - Care, Cultivation, and Display}}
       }{Stefano Bragaglia}
\label{2025-04-mini}

\begin{figure*}[!b]
\centering
\includegraphics[width=0.85\textwidth]{2025_04/res/intro}
\centerline{\emph{Historic Chobham Village Hall welcomed bonsai lovers of \textbf{\textit{all}} kinds.}}
\end{figure*}

\noindent
At the recent UK Bonsai Association workshop on Mini Bonsai on March 29, 2025, held at Chobham Village Hall, we had the pleasure of hearing from some fantastic speakers, each sharing their expertise on different aspects of mame bonsai.
Mark Moreland kicked things off with an insightful talk on the importance of humidity in bonsai care, while Alex Rudd delved into the world of pots, discussing the best choices for mame.
David Cheshire then introduced us to traditional Japanese growing techniques for these miniature trees, and Gary Careford rounded off the session with an engaging discussion on the principles of displaying mame bonsai effectively.

Alongside the four presentations, the event boasted four packed rows of shohin and mame bonsai displays, each brimming with delightful little trees, exquisite stands, and thoughtfully curated accents and scrolls—not to mention plenty of clever touches to highlight the plants.
The contributors included some big names in the UK bonsai scene, like Mark and Mingchen Moreland and Gary Careford, as well as plenty of passionate enthusiasts.
In the afternoon, the exhibition opened to the public and drew a healthy crowd.
To round things off, there were a handful of vendor tables offering everything from pots (including Alex Rudd’s fine selection) and trees (with David Cheshire among them) to stands, scrolls, and all sorts of other must-have accessories for mini bonsai lovers.

In the following articles, we'll be covering the key takeaways from each of these talks, offering valuable insights for those who couldn't attend and a refresher for those who did. % Whether you're new to mame bonsai or a seasoned enthusiast, there's plenty to learn from these experts!
